% !TeX program = xelatex
% 14-page editable starter following the supplied PPTX's page sequence.
\documentclass[10pt,aspectratio=169]{beamer}
\usetheme{PCL} % 如需额外章节过渡页，可用 \usetheme[sectionpages]{PCL}
\usepackage{amsmath,amssymb,booktabs,tabularx}
\title[报告简称]{题目：在此填写报告标题}
\subtitle{}
\author{汇报人：XXX}
\institute{XXX研究部}
\date{\PCLtoday}
\contact{}

\begin{document}
\maketitle
\PCLcontents

\section{研究背景}
\begin{frame}{一、研究背景}
  \PCLtakeaway[研究问题]{在此填写项目面向的问题与需求。}
  \vspace{0.5cm}
  \begin{itemize}
    \setlength{\itemsep}{0.7em}
    \item 研究背景
    \item 现有研究与主要不足
    \item 开展本项目的必要性
  \end{itemize}
\end{frame}

\section{研究目标}
\begin{frame}{二、研究目标}
  \begin{block}{总体目标}
    在此填写项目总体目标。
  \end{block}
  \begin{enumerate}
    \setlength{\itemsep}{0.7em}
    \item 目标一
    \item 目标二
    \item 目标三
  \end{enumerate}
\end{frame}

\section{研究内容}
\begin{frame}{三、研究内容}
  \begin{itemize}
    \setlength{\itemsep}{0.8em}
    \item 研究内容一
    \item 研究内容二
    \item 研究内容三
  \end{itemize}
\end{frame}

\begin{frame}{项目分解}
  \centering
  \PCLplaceholder[3.8cm]{12cm}{在此插入课题分解图}
\end{frame}

\begin{frame}{课题 1：XXX研究内容}
  \begin{columns}[T,onlytextwidth]
    \begin{column}{0.47\textwidth}
      \begin{block}{研究任务}
        在此填写课题的关键研究任务。
      \end{block}
    \end{column}
    \begin{column}{0.47\textwidth}
      \begin{block}{技术路线}
        在此填写方法与技术路线。
      \end{block}
    \end{column}
  \end{columns}
\end{frame}

\begin{frame}{课题 1：预期成果}
  \begin{itemize}
    \setlength{\itemsep}{0.8em}
    \item 成果内容
    \item 考核指标
    \item 交付形式
  \end{itemize}
\end{frame}

\begin{frame}{项目实施计划}
  \centering
  \PCLplaceholder[3.8cm]{12cm}{在此填写阶段计划与里程碑}
\end{frame}

\section{预期成果}
\begin{frame}{四、项目预期代表性成果}
  \begin{enumerate}
    \setlength{\itemsep}{0.8em}
    \item 代表性成果一
    \item 代表性成果二
    \item 代表性成果三
  \end{enumerate}
\end{frame}

\begin{frame}{项目特色与关键突破}
  \PCLtakeaway[关键特色]{在此填写本项目的核心差异及支撑依据。}
\end{frame}

\section{研究团队}
\begin{frame}{五、研究团队}
  \renewcommand{\arraystretch}{1.7}
  \begin{tabularx}{\linewidth}{@{}llX@{}}
    \toprule
    角色 & 姓名 / 团队 & 主要职责\\
    \midrule
    项目负责人 & XXX & 待填写\\
    课题负责人 & XXX & 待填写\\
    研究骨干 & XXX & 待填写\\
    \bottomrule
  \end{tabularx}
\end{frame}

\begin{frame}{研究基础}
  \centering
  % 使用实际图片：\includegraphics[width=0.9\linewidth,height=4cm,keepaspectratio]{assets/your-figure.pdf}
  \PCLplaceholder[3.8cm]{12cm}{在此展示前期工作与研究条件}
\end{frame}

\PCLthanks
\end{document}
